\def\level{Première}  % Classe à droite
\def\course{NSI}
\def\chaptername{Tris par insertion et sélection} % Titre au centre
\def\docname{AP08-S}        % Identifiant à gauche

\pagestyle{cours_FP}
\makeheader{} % Affiche le bandeau

	\begin{minipage}{0.48\linewidth}
\begin{propriete}[Le tri par insertion]{}
		\includegraphics[width=0.95\linewidth]{tri-par-insertion.png}
	\end{propriete}

	\end{minipage}\hfill{}
	\begin{minipage}{0.48\linewidth}
\begin{propriete}[Le tri par sélection]{}
		\includegraphics[width=0.95\linewidth]{tri-par-selection.png}
	\end{propriete}

	\end{minipage}



\begin{propriete}[Le tri par insertion]{}
\begin{CodePythontexAlt}[TaillePolice=\large, Largeur=1\linewidth,Centre,PremLigne=1]{}
		def tri_insertion(L):
			''' Algorithme de tri par insertion '''
			n = len(L)
			for i in range(1, n):
				# Indice de la clé
				k = i				
				# Tant que le bord de la liste n'est pas atteint
				# Et que la clé est inférieure à l'élement de gauche:
				while k > 0 and L[k] < L[k-1]:
					# Échange de la clé avec l'élément de gauche
					L[k], L[k-1] = L[k-1], L[k]
					# Décalage de la clé vers la gauche
					k -= 1
		\end{CodePythontexAlt}
\end{propriete}


	
\begin{propriete}[Le tri par sélection]{}
\begin{CodePythontexAlt}[TaillePolice=\large, Largeur=1\linewidth,Centre,PremLigne=1]{}
		def tri_selection(L):
			''' Algorithme de tri par sélection '''
			n = len(L)
			for i in range(0, n-1):
				# Recherche du plus petit element de la liste non triee
				mini = i		# indice de plus petit element non triee
				for j in range(i+1, n):
					if L[j] < L[mini]:
						mini = j 		# j est le nouvel indice du minimum	
				# Echange entre L[i] et L[mini]
				L[i], L[mini] = L[mini], L[i]
		\end{CodePythontexAlt}

	


\end{propriete}



	
